% problem-000463 4 煤气化制氢热力学
\section*{4. 煤气化制氢热力学}
氢最有可能成为 21 世纪的主要能源，但氢⽓需要由其他物质来制备，制氢的⽅法之⼀是以煤的转化为基础。基本原理是⽤碳、⽔在⽓化炉中发⽣如下反应：

$
\mathrm{C} (\mathrm{s}) + \mathrm{H} _ {2} \mathrm{O} \rightleftharpoons \mathrm{CO} (\mathrm{g}) + \mathrm{H} _ {2} (\mathrm{g})\tag{1}
$

$
\mathrm{CO} (\mathrm{g}) + \mathrm{H} _ {2} \mathrm{O} \rightleftharpoons \mathrm{CO} _ {2} (\mathrm{g}) + \mathrm{H} _ {2} (\mathrm{g})\tag{2}
$

利⽤CaO吸收产物中的CO_{2}：

$
\mathrm{CaO} (\mathrm{s}) + \mathrm{CO} _ {2} (\mathrm{g}) \rightleftharpoons \mathrm{CaCO} _ {3} (\mathrm{s})\tag{3}
$

产物中的 $\mathrm { H } _ { 2 }$ 与平衡体系中的 $\mathrm { C } ,$ ，CO， $\mathrm { C O _ { 2 } }$ 发⽣反应，⽣成CH_{4}：

$
\mathrm{C} (\mathrm{s}) + 2 \mathrm{H} _ {2} (\mathrm{g}) \rightleftharpoons \mathrm{CH} _ {4} (\mathrm{g})\tag{4}
$

$
\mathrm{CO} (\mathrm{g}) + 3 \mathrm{H} _ {2} (\mathrm{g}) \rightleftharpoons \mathrm{CH} _ {4} (\mathrm{g}) + \mathrm{H} _ {2} \mathrm{O} (\mathrm{g})\tag{5}
$

$
\mathrm{CO} _ {2} (\mathrm{g}) + 4 \mathrm{H} _ {2} (\mathrm{g}) \rightleftharpoons \mathrm{CH} _ {4} (\mathrm{g}) + 2 \mathrm{H} _ {2} \mathrm{O} (\mathrm{g})\tag{6}
$

将 2 mol C(s)，2 mol $\mathrm { H _ { 2 } O ( g ) }$ ，2 mol $\mathrm { C a O } ( s )$ 放⼊⽓化炉，在 $8 5 0 ~ ^ { \circ } \mathrm { C }$ 下发⽣反应。

已知850 °C下相关物种的热⼒学参数：

\textit{（原卷此处含表格/仪器试剂表，详见 PDF 或 MD 源文件。）}

\subsection*{4-1}
计算⽓化炉总压为 $2 . 5 0 \times 1 0 ^ { 6 } \mathrm { P a }$ 时， $\mathrm { H } _ { 2 }$ 在平衡混合⽓中的摩尔分数。

\subsection*{4-2}
计算 $8 5 0 ^ { \circ } \mathrm { C }$ 从起始原料到平衡产物这⼀过程的热效应。

\subsection*{4-3}
碳在⾼温下是⼀种优良的还原剂，可⽤于冶炼多种⾦属。试写出 $6 0 0 ^ { \circ } \mathrm { C }$ 碳的可能氧化产物的化学式。从热⼒学⻆度说明原因（假设 $6 0 0 ^ { \circ } \mathrm { C }$ 反应的熵变、焓变和 $8 5 0 ^ { \circ } \mathrm { C }$ 下的熵变、焓变相同）
